\documentclass[10pt,a4paper]{amsart}
\usepackage{amsmath,amssymb,amsthm,graphicx}
% Retained minimal preamble for the Stacks Project chapter more-algebra.tex (More on Algebra).
% The chapter file pinned-original-source.tex carries only \input{preamble} on its first line, so no
% preamble text lives inside the chapter itself.  The Stacks Project uses one single shared preamble for
% every chapter; the snapshot stages it identically as sources/stacks-chow/preamble.tex and
% sources/stacks-etale-cohomology/preamble.tex.  The block below is copied mechanically, line for line,
% from sources/stacks-chow/preamble.tex lines 14-18 (the Xy-pic setup) and 154-217 (the native theorem
% declarations together with the mathematical macros).  Nothing was added, removed or redefined; the
% wrapper class and the amsmath/amssymb/amsthm/graphicx base packages are supplied by the assembler.
% Chapter: more-algebra.tex -- More on Algebra.
% For commutative diagrams we use Xy-pic
\usepackage[all]{xy}

% We use 2cell for 2-commutative diagrams.
\xyoption{2cell}

% Theorem environments.
%
\theoremstyle{plain}
\newtheorem{theorem}[subsection]{Theorem}
\newtheorem{proposition}[subsection]{Proposition}
\newtheorem{lemma}[subsection]{Lemma}

\theoremstyle{definition}
\newtheorem{definition}[subsection]{Definition}
\newtheorem{example}[subsection]{Example}
\newtheorem{exercise}[subsection]{Exercise}
\newtheorem{situation}[subsection]{Situation}

\theoremstyle{remark}
\newtheorem{remark}[subsection]{Remark}
\newtheorem{remarks}[subsection]{Remarks}

\numberwithin{equation}{subsection}

% Macros
%
\def\lim{\mathop{\mathrm{lim}}\nolimits}
\def\colim{\mathop{\mathrm{colim}}\nolimits}
\def\Spec{\mathop{\mathrm{Spec}}}
\def\Hom{\mathop{\mathrm{Hom}}\nolimits}
\def\Ext{\mathop{\mathrm{Ext}}\nolimits}
\def\SheafHom{\mathop{\mathcal{H}\!\mathit{om}}\nolimits}
\def\SheafExt{\mathop{\mathcal{E}\!\mathit{xt}}\nolimits}
\def\Sch{\mathit{Sch}}
\def\Mor{\mathop{\mathrm{Mor}}\nolimits}
\def\Ob{\mathop{\mathrm{Ob}}\nolimits}
\def\Sh{\mathop{\mathit{Sh}}\nolimits}
\def\NL{\mathop{N\!L}\nolimits}
\def\CH{\mathop{\mathrm{CH}}\nolimits}
\def\proetale{{pro\text{-}\acute{e}tale}}
\def\etale{{\acute{e}tale}}
\def\QCoh{\mathit{QCoh}}
\def\Ker{\mathop{\mathrm{Ker}}}
\def\Im{\mathop{\mathrm{Im}}}
\def\Coker{\mathop{\mathrm{Coker}}}
\def\Coim{\mathop{\mathrm{Coim}}}

% Boxtimes
%
\DeclareMathSymbol{\boxtimes}{\mathbin}{AMSa}{"02}

%
% Macros for moduli stacks/spaces
%
\def\QCohstack{\mathcal{QC}\!\mathit{oh}}
\def\Cohstack{\mathcal{C}\!\mathit{oh}}
\def\Spacesstack{\mathcal{S}\!\mathit{paces}}
\def\Quotfunctor{\mathrm{Quot}}
\def\Hilbfunctor{\mathrm{Hilb}}
\def\Curvesstack{\mathcal{C}\!\mathit{urves}}
\def\Polarizedstack{\mathcal{P}\!\mathit{olarized}}
\def\Complexesstack{\mathcal{C}\!\mathit{omplexes}}
% \Pic is the operator that assigns to X its picard group, usage \Pic(X)
% \Picardstack_{X/B} denotes the Picard stack of X over B
% \Picardfunctor_{X/B} denotes the Picard functor of X over B
\def\Pic{\mathop{\mathrm{Pic}}\nolimits}
\def\Picardstack{\mathcal{P}\!\mathit{ic}}
\def\Picardfunctor{\mathrm{Pic}}
\def\Deformationcategory{\mathcal{D}\!\mathit{ef}}
\usepackage{cleveref}
\crefname{originalequation}{original equation}{original equations}
\Crefname{originalequation}{Original equation}{Original equations}
\makeatletter
\newlabel{equation-BL-cech-mod-re}{{0F1R}{source}{}{sourcepointer-1}{}}
\newlabel{equation-BL-cech-mod-re@cref}{{[originalequation][0][]0F1R}{[1][0][]source}{}{}{}}
\makeatother
\crefname{originallemma}{original lemma}{original lemmas}
\Crefname{originallemma}{Original lemma}{Original lemmas}
\makeatletter
\newlabel{algebra-lemma-tensor-product-exact}{{00DF}{source}{}{sourcepointer-2}{}}
\newlabel{algebra-lemma-tensor-product-exact@cref}{{[originallemma][0][]00DF}{[1][0][]source}{}{}{}}
\makeatother
\crefname{originallemma}{original lemma}{original lemmas}
\Crefname{originallemma}{Original lemma}{Original lemmas}
\makeatletter
\newlabel{lemma-same-f-torsion}{{0BNR}{source}{}{sourcepointer-3}{}}
\newlabel{lemma-same-f-torsion@cref}{{[originallemma][0][]0BNR}{[1][0][]source}{}{}{}}
\makeatother
\title{Exactness of the Cech sequence of a module is equivalent to bijectivity of the comparison map on f-power torsion}
\author{The Stacks Project Authors}
\date{stacks-project:more-algebra.tex:local-snapshot-20261012; GNU FDL 1.2 or later}
\begin{document}
\maketitle
\setcounter{subsection}{1}
\noindent\textit{Editorial note.} All delivered mathematical text is the author's own, copied from the pinned Stacks Project chapter file more-algebra.tex; no statement or proof step was rewritten, reordered or removed. The counted statement is the labeled lemma environment at source lines 25364-25385 (canonical tag 0BNW) and the counted proof is the proof environment at source lines 25387-25423, which immediately follows it and is its own proof.\par\medskip
\noindent\textit{Editorial note.} The statement refers to the displayed Cech sequence by its label equation-BL-cech-mod-re, which lies outside the counted spans; that label is carried as a navigation stub so that the reference resolves in the compiled document.\par\medskip
\noindent\textit{Editorial note.} The two other references are to right exactness of the tensor product in another chapter and to the sibling lemma of this chapter; both are navigation stubs with their roles recorded in the reasons.\par\medskip
\noindent\textit{Editorial note.} The chapter file pinned-original-source.tex carries only an input of the shared preamble on its first line, so no preamble text lives inside the chapter. The retained file sources/stacks-more-algebra/retained-preamble.tex is a mechanical, line-for-line copy of the project's single shared preamble as staged in the snapshot (Xy-pic setup plus the native theorem declarations and mathematical macros). Nothing was added, removed or redefined, and the excerpt serves only as preamble, never as a span coordinate.\par\medskip
\noindent\textit{Editorial note.} The source URL is the source-level Stacks Project URL, as required when more than one card is produced from one source; the per-claim tag is carried by claim\_key and by this note. The counted statement and proof cite no bibliography entry, so this card needs no bibliography.\par\medskip
\noindent\textit{Allowed standard prerequisites (editorial).}\begin{itemize}
\item The Cech sequence of a module for the pair consisting of a ring map and an element compares the module with its tensor product with the target ring and its localization. used at proof lines 25416-25419, where the sequence labelled equation-BL-cech-mod-re is referred to. Fixes the exact sequence whose exactness in each position is characterized by the lemma.
\item Tensoring with a module is right exact, so the sequence in question is always exact on the right. invoked at proof lines 25389-25391. Gives part (1) immediately.
\item The kernel of the localization map is the submodule of elements annihilated by a power of the element, which identifies the left exactness condition. used at proof lines 25393-25395. Gives part (2).
\end{itemize}\medskip
\ifdefined\hypertarget\hypertarget{sourcepointer-1}{}\fi
\noindent\textit{Source reference.} Original equation 0F1R: sources/stacks-more-algebra/pinned-original-source.tex, source lines 25346--25346. The reference points to the displayed Cech sequence that the statement and proof are about; it names the object under discussion and the element computations that decide exactness in each position are inside the counted span.\par\medskip
\ifdefined\hypertarget\hypertarget{sourcepointer-2}{}\fi
\noindent\textit{Source reference.} Original lemma 00DF: sources/stacks-chow/tags.txt, source lines 495--495. The reference records right exactness of the tensor product, invoked at proof lines 25389-25391 for part (1) and again at proof lines 25416-25418 for the exactness in the middle for glueing pairs; it is a standard pointer and everything drawn from it is inside the counted span.\par\medskip
\ifdefined\hypertarget\hypertarget{sourcepointer-3}{}\fi
\noindent\textit{Source reference.} Original lemma 0BNR: sources/stacks-more-algebra/pinned-original-source.tex, source lines 25255--25255. The reference identifies the tensor product of the ring modulo its power torsion with the target modulo its own power torsion, used at proof lines 25388-25389 and 25399-25401; the exactness computation built on it is inside the counted span.\par\medskip
\begin{lemma}
\label{lemma-same-f-torsion-module}
Let $R$ be a ring, let $f \in R$, and let $R \to R'$ be a ring map
which induces isomorphisms $R/f^nR \to R'/f^nR'$ for $n > 0$.
The sequence (\ref{equation-BL-cech-mod-re}) is
\begin{enumerate}
\item exact on the right,
\item exact on the left if and only if
$M[f^\infty] \to (M \otimes_R R')[f^\infty]$
is injective, and
\item exact in the middle if and only if
$M[f^\infty] \to (M \otimes_R R')[f^\infty]$
is surjective.
\end{enumerate}
Thus $M$ is glueable for $(R \to R', f)$ if and only if
$M[f^\infty] \to (M \otimes_R R')[f^\infty]$ is bijective.
If $(R \to R', f)$ is a glueing pair, then $M$ is glueable for $(R \to R', f)$
if and only if $M[f^\infty] \to (M \otimes_R R')[f^\infty]$ is injective.
For example, if $(R, f)$ is a glueing pair, then $M$ is glueable
if and only if $M[f^\infty] \to (M \otimes_R R^\wedge)[f^\infty]$ is
injective.
\end{lemma}

\begin{proof}
We will use the results of Lemma \ref{lemma-same-f-torsion} without
further mention. The functor $M \otimes_R -$ is right exact
(Algebra, Lemma \ref{algebra-lemma-tensor-product-exact})
hence we get (1).

\medskip\noindent
The kernel of $M \to M \otimes_R R_f = M_f$ is $M[f^\infty]$.
Thus (2) follows.

\medskip\noindent
If the sequence is exact in the middle, then elements of the form
$(x, 0)$ with $x \in (M \otimes_R R')[f^\infty]$ are in the image of
the first arrow. This implies that
$M[f^\infty] \to (M \otimes_R R')[f^\infty]$
is surjective. Conversely, assume that
$M[f^\infty] \to (M \otimes_R R')[f^\infty]$
is surjective. Let $(x, y)$ be an element in the middle
which maps to zero on the right. Write $y = y'/f^n$ for some $y' \in M$.
Then we see that $f^n x - y'$ is annihilated by some power of $f$ in
$M \otimes_R R'$. By assumption we can write $f^nx - y' = z$ for some
$z \in M[f^\infty]$. Then $y = y''/f^n$ where $y'' = y' + z$
is in the kernel of $M \to M/f^nM$. Hence we see that $y$ can be
represented as $y'''/1$ for some $y''' \in M$. Then
$x - y'''$ is in $(M \otimes_R R')[f^\infty]$.
Thus $x - y''' = z' \in M[f^\infty]$.
Then $(x, y'''/1) = (y''' + z', (y''' + z')/1)$ as desired.

\medskip\noindent
If $(R \to R', f)$ is a glueing pair, then (\ref{equation-BL-cech-mod-re})
is exact in the middle for any $M$ by
Algebra, Lemma \ref{algebra-lemma-tensor-product-exact}.
This gives the penultimate statement of the lemma.
The final statement of the lemma follows from this and
the fact that $(R, f)$ is a glueing pair if and only if
$(R \to R^\wedge, f)$ is a glueing pair.
\end{proof}

\end{document}
